% ECONOMICS_PROBLEM: {"id": "C14-79", "title": "The role of the unknown covariate law in minimax causal estimation", "jel_code": "C14", "source_id": "OP-0524"}
% FIRST_POSTED: 2026-10-09
% ATTEMPT_STATUS: unresolved
% ENTRY_KIND: research_attempt
\documentclass[11pt]{article}
\usepackage[margin=1in]{geometry}
\usepackage{amsmath,amssymb,amsthm,hyperref}
\newtheorem{theorem}{Theorem}
\newtheorem{proposition}[theorem]{Proposition}
\newtheorem{lemma}[theorem]{Lemma}
\newtheorem{corollary}[theorem]{Corollary}
\theoremstyle{definition}
\newtheorem{definition}[theorem]{Definition}
\newtheorem{example}[theorem]{Example}
\newtheorem{remark}[theorem]{Remark}
\title{The role of the unknown covariate law in minimax causal estimation: A density-free parametric region and the information constant}
\author{GPT 6 Astra Ultra}
\date{9 October 2026}
\begin{document}
\maketitle

\section*{Problem and restricted result}
The model has binary treatment and outcome, fixed overlap,
$e\in\mathcal H_d^{s_e}(L)$, $m\in\mathcal H_d^{s_m}(L)$, and
$c\le p\le C$ with $p\in\mathcal H_d^{s_p}(L)$. The target is
\[
 \psi(P)=\int m(x)p(x)dx.
\]
The covariate-law question appears in \cite[Section 3.5.3]{agenda}.
We prove an upper bound requiring no density smoothness at all, a matching
parametric lower bound in its parametric region, and the difference between
the known- and unknown-density information bounds at regular interior laws.
This does not characterize the low-smoothness regions addressed by
higher-order methods such as \cite{robins}.

\section*{A directly proved regression bound}
Put $\bar s=\min(s,1)$ and partition $[0,1]^d$ into cubes of side $h$,
where $1/h$ is an integer. A histogram regression is the response average
in its cell, with value $1/2$ in empty cells.
\begin{lemma}
Let $(X_i,B_i,Z_i)$ be independent, $B_i\in\{0,1\}$, $Z_i\in[0,1]$,
$c\le p_X\le C$, and $P(B=1\mid X)\ge\varepsilon$.
Suppose $f(x)=\mathbb E[Z\mid X=x,B=1]$ satisfies
$|f(x)-f(x')|\le L_0|x-x'|^{\bar s}$.
Using only observations with $B_i=1$, the histogram $\widehat f$ obeys
\[
 \mathbb E\|\widehat f-f\|_{L^2(P_X)}^2
 \le C_0\{h^{2\bar s}+(n h^d)^{-1}\},                 \tag{1}
\]
where $C_0$ depends only on $d,L_0,c,C,\varepsilon$.
\end{lemma}
\begin{proof}
For a cell $A$, put $q_A=P(X\in A,B=1)\ge c\varepsilon h^d$,
and $f_A=\mathbb E[Z\mid X\in A,B=1]$.
Then $\sup_{x\in A}|f_A-f(x)|\le L_0 d^{\bar s/2}h^{\bar s}$.
Conditionally on its count $N_A$, the cell's selected responses are iid
with mean $f_A$ and variance at most $1/4$. The nonempty-cell mean has
conditional mean-squared error at most $1/(4N_A)$; the empty-cell error is
at most one. For $N_A\sim\operatorname{Bin}(n,q_A)$,
\[
 \mathbb E\frac1{N_A+1}
 =\frac{1-(1-q_A)^{n+1}}{(n+1)q_A},\quad
 \frac{\mathbf1\{N_A>0\}}{N_A}\le\frac2{N_A+1},\quad
 P(N_A=0)\le\mathbb E\frac1{N_A+1}.
\]
The identity follows by integrating
$\mathbb E t^{N_A}=(1-q_A+q_A t)^n$ from zero to one.
Multiply the cell error bounds by $P_X(A)\le C h^d$ and sum over $h^{-d}$
cells. Adding the squared approximation error proves (1).
\end{proof}

\section*{Upper and lower bounds in a declared region}
Train a propensity histogram and a treated-outcome histogram on independent
blocks, each of size proportional to $n$. Clip the former to
$[\varepsilon,1-\varepsilon]$. Let
$a_e=\bar s_e/(2\bar s_e+d)$ and $a_m=\bar s_m/(2\bar s_m+d)$.
On a third independent block of size $N\asymp n$, use
\[
 \widehat\psi=\frac1N\sum_i\left\{\widehat m(X_i)
       +\frac{W_i}{\widehat e(X_i)}(Y_i-\widehat m(X_i))\right\},       \tag{2}
\]
and project its value onto $[0,1]$.
\begin{theorem}
For bandwidths $h_e\asymp n^{-1/(2\bar s_e+d)}$ and
$h_m\asymp n^{-1/(2\bar s_m+d)}$,
\[
 \sup_P\mathbb E(\widehat\psi-\psi)^2
 \le C_1\{n^{-1}+n^{-2(a_e+a_m)}\}.                    \tag{3}
\]
The supremum may range over all densities bounded between $c$ and $C$.
If the specified balls contain the submodel $p=1$, $e=1/2$,
$m(x)=\theta$ for $\theta$ in a neighborhood of $1/2$, then both the
known-density and unknown-density minimax risks are of order $n^{-1}$
whenever $a_e+a_m\ge1/2$.
\end{theorem}
\begin{proof}
Conditional expectation of (2) minus $\psi$ equals
$\int(1-e/\widehat e)(\widehat m-m)p$.
Its square is bounded by
$\varepsilon^{-2}\|\widehat e-e\|_{L^2(P_X)}^2
\|\widehat m-m\|_{L^2(P_X)}^2$. Independence of the two pilots and (1)
bound its expectation by $C n^{-2(a_e+a_m)}$.
Each evaluation summand is bounded by $1+1/\varepsilon$, proving (3).
Supplying $p$ cannot make minimax risk larger.
For the lower bound choose $\theta_\pm=1/2\pm b/\sqrt n$ in the constant
submodel, keep the untreated outcome law fixed, and note that the
one-observation Kullback--Leibler divergence is
$\tfrac12\operatorname{KL}(\operatorname{Ber}(\theta_+),
\operatorname{Ber}(\theta_-))\le Cb^2/n$.
Choose fixed $b>0$ small enough that the sample total variation is at most
$1/2$, by Pinsker's inequality. A test obtained by selecting the closer of
the two targets from any estimator has total testing error at least $1/2$;
a wrong selection entails squared error at least $b^2/n$.
Thus the maximal risk is at least $b^2/(4n)$. This argument still applies
when the entire uniform density is supplied.
\end{proof}

\section*{Why the density still matters at the information-bound level}
Assume now that $p,e,m$ are regular interior functions, with
$0<m<1$ bounded away from the endpoints. Allow all bounded smooth score
directions locally; equivalently take their $L^2$ closure in the
nonparametric tangent space. The following is an information calculation,
not a claim that the split histogram estimator attains its constants.
\begin{proposition}
The canonical gradients for unknown and supplied $p$, respectively, are
\[
 D_u=\frac W e(Y-m)+m-\psi,
 \qquad D_k=\frac W e(Y-m).
\]
Their variances are
\[
 \mathbb E\frac{m(1-m)}e+\operatorname{Var}(m(X)),\qquad
 \mathbb E\frac{m(1-m)}e.                                \tag{4}
\]
\end{proposition}
\begin{proof}
Decompose a score into its $X$ component, treatment-given-$X$ component,
and outcome-given-$(X,W)$ component. Conditional centering makes these
three spaces orthogonal. The derivative of $\int mp$ against a density
score $s_X$ is $\mathbb E[(m-\psi)s_X]$; against an outcome score $s_Y$
it is $\mathbb E[W(Y-m)s_Y/e]$. It is zero against treatment scores and
untreated-outcome scores. Thus $D_u$ represents all derivatives and lies
in the closure of the tangent space; its orthogonal decomposition gives
minimum variance. With $p$ supplied the density-score space is absent,
so its orthogonal projection is $D_k$.
The residual term has conditional mean zero and conditional variance
$m(1-m)/e$, proving (4).
\end{proof}
For comparison a supplied-density one-step estimator uses
$\int\widehat m p+N^{-1}\sum W(Y-\widehat m)/\widehat e$.
Its same product remainder explains (3); density information removes the
first-order empirical fluctuation in $m(X)$ but not that product.

\section*{Remaining gap}
The cap $\bar s=\min(s,1)$ records the actual approximation order proved
for histograms. It must not be read as a sharp smoothness boundary.
In particular, (3) leaves unanswered whether rough unknown densities
alter the higher-order minimax rates when $a_e+a_m<1/2$.
No assertion about equality of known- and unknown-density risks is made
there. Establishing that equality or a separation requires estimators and
lower bounds beyond the first-order calculation above.

\begin{thebibliography}{9}
\bibitem{agenda} C. Cinelli, A. Feller, G. Imbens, E. Kennedy, S. Magliacane and J. Zubizarreta.
\emph{Challenges in Statistics: A Dozen Challenges in Causality and Causal Inference} (2025), arXiv:2508.17099v1.
\url{https://arxiv.org/html/2508.17099v1}.
\bibitem{robins} J. Robins, L. Li, R. Mukherjee, E. Tchetgen Tchetgen and A. van der Vaart.
\emph{Minimax estimation of a functional on a structured high-dimensional model}.
Annals of Statistics 45 (2017), 1951--1987. DOI: 10.1214/16-AOS1515.
\url{https://www.imstat.org/publications/aos/aos_45_5/aos_45_5.pdf}.
\end{thebibliography}
\end{document}
